\section{Inverse design is a settling least-fixed-point}

\begin{theorem}[Inverse design is a least-fixed-point over the carrier]\label{thm:inverse-design-least-fixed-point}\notready\uses{thm:thickened-torus-z2-carrier, thm:homogenized-energy-settled}
\emph{(Open frontier --- stated, not formalized.)} Inverse design is a settling
problem on the $\mathbb{Z}^2$ carrier
(\cref{thm:thickened-torus-z2-carrier}): the design objective
$\mathcal{L}(s)=d(f(X),M)+\lambda_1\operatorname{Err}_{\text{pat}}+\lambda_2\operatorname{Err}_{\text{mech}}$
is minimized at the least fixed point of the local-constraint settle, and the settled
configuration is the manufacturable programme. The homogenized energy
(\cref{thm:homogenized-energy-settled}) is the mechanical term's settled value. Forward
weaving is the propagator; inverse design is its settling inverse. Requires the design
functional as a monotone operator on the carrier.
\end{theorem}

\begin{theorem}[Manufacturability is a settling predicate]\label{thm:manufacturability-settling-predicate}\notready\uses{thm:inverse-design-least-fixed-point, thm:bead-threading-closed-walk}
\emph{(Open frontier --- stated, not formalized.)} The predicate
$\mathrm{Manufacturable}(p)=\mathrm{true}$ says that the programme's
local-compatibility system settles without a held
$\top$-defect: a draft is manufacturable iff its neighbourhood constraint reaches a fixed
point (no float, lace, or port violation), and a beadwork is manufacturable iff its threading
walk closes (\cref{thm:bead-threading-closed-walk}). Manufacturability is settlement
reachability; the verified textile compiler is the decision procedure for that fixed point.
Requires the manufacturability predicate as settle-reachability on the carrier.
\end{theorem}
